% ══════════════════════════════════════════════════════════════════
% Section 4 — Training Procedure
% ══════════════════════════════════════════════════════════════════
\section{Training procedure}
\label{sec:training}

\subsection{Data generation}
\label{sec:training:data}

Training data are generated synthetically by sampling deformation gradients
$\Fmat$ and computing the analytical constitutive response.

\subsubsection{Deformation sampling}

The \texttt{DeformationGenerator} produces deformation gradients via three
modes combined with random rotations for frame invariance:
\begin{itemize}
  \item \textbf{Uniaxial:} $\Fmat = \text{diag}(\lambda, \lambda^{-1/2},
    \lambda^{-1/2})$, with $\lambda \in [\lambda_\min, \lambda_\max]$.
  \item \textbf{Biaxial:} $\Fmat = \text{diag}(\lambda_1, \lambda_2,
    (\lambda_1\lambda_2)^{-1})$ (volume-preserving).
  \item \textbf{Shear:} $F_{12} = \gamma$, $\gamma \in [-\gamma_\max,
    \gamma_\max]$.
\end{itemize}
A composite generator (\texttt{combined}) draws equal proportions from all
three modes and applies random rotations $\tens{Q}$:
$\Fmat \leftarrow \tens{Q}\Fmat$.

For isotropic materials, the stretch range is $[0.7, 1.5]$ with shear up to
$0.3$.  Anisotropic materials use a tighter range $[0.85, 1.2]$ with shear up
to $0.15$ to avoid exponential overflow in fiber terms.

\subsubsection{Input representation}

Invariants are computed from the sampled deformation gradients:
\begin{equation}
  \tens{x}_\text{iso} = (\bar{I}_1, \bar{I}_2, J) \in \R^3, \qquad
  \tens{x}_\text{aniso} = (\bar{I}_1, \bar{I}_2, J, I_4, I_5) \in \R^5.
  \label{eq:input_repr}
\end{equation}
Inputs are standardized to zero mean and unit variance using a fitted
normalizer: $\hat{\tens{x}} = (\tens{x} - \bm{\mu}) / \bm{\sigma}$.

\subsubsection{Target computation}

For each sample, the analytical material model provides:
\begin{itemize}
  \item Strain energy: $W = \sef(\Cmat)$.
  \item Invariant gradients: $\partial \sef / \partial I_\alpha$, computed via
    symbolic differentiation (SymPy) for anisotropic materials or via
    analytical expressions for isotropic materials.
\end{itemize}
The invariant gradients are rescaled by the input standard deviation:
$\hat{S}_\alpha = (\partial\sef/\partial I_\alpha) \cdot \sigma_\alpha$ to
match the chain-rule relationship with the normalized inputs.

\subsection{Energy--stress loss}
\label{sec:training:loss}

Standard regression on energy alone produces models with poor stress accuracy
because small energy errors can translate to large gradient errors.  Conversely,
training on stress alone discards the potential structure.  We adopt a dual
loss that penalizes both:
\begin{equation}
  \mathcal{L} = \alpha\,\mathcal{L}_W + \beta\,\mathcal{L}_S,
  \label{eq:energy_stress_loss}
\end{equation}
where
\begin{align}
  \mathcal{L}_W &= \frac{1}{N}\sum_{i=1}^{N}
    \bigl(\hat{\sef}(\hat{\tens{x}}_i) - W_i\bigr)^2, \label{eq:loss_energy} \\
  \mathcal{L}_S &= \frac{1}{Nd}\sum_{i=1}^{N}\sum_{\alpha=1}^{d}
    \biggl(\pd{\hat{\sef}}{\hat{x}_\alpha}\bigg|_{\hat{\tens{x}}_i}
    - \hat{S}_{i\alpha}\biggr)^2, \label{eq:loss_stress}
\end{align}
and $\partial\hat{\sef}/\partial\hat{x}_\alpha$ is computed via PyTorch
automatic differentiation through the network.  We use $\alpha = \beta = 1$
throughout.

This formulation is architecture-agnostic: it applies identically to MLP,
ICNN, and polyconvex ICNN, and is essential for the convex architectures
(which output scalar energy only).

\subsection{Training details}
\label{sec:training:details}

\begin{table}[htbp]
\centering
\caption{Training hyperparameters.}
\label{tab:training_hyperparams}
\begin{tabular}{@{}ll@{}}
\toprule
\textbf{Parameter} & \textbf{Value} \\
\midrule
Optimizer          & Adam ($\text{lr} = 10^{-3}$) \\
Learning rate schedule & ReduceLROnPlateau (factor 0.5, patience 20) \\
Batch size         & 256 \\
Maximum epochs     & 1000 \\
Early stopping     & patience = 200 (validation loss) \\
Train / val split  & 85\% / 15\% \\
Samples (isotropic)   & 10\,000 \\
Samples (anisotropic) & 10\,000 \\
Loss weights       & $\alpha = 1$, $\beta = 1$ \\
Activation         & softplus (all architectures) \\
\bottomrule
\end{tabular}
\end{table}

The training pipeline is implemented in the \texttt{Trainer} class, which
handles batching, early stopping, learning rate scheduling, and history
logging.  All experiments use a single random seed for reproducibility.

% ── Training pipeline diagram ─────────────────────────────────────
\begin{figure}[htbp]
\centering
% TODO: Replace with a flowchart showing:
%   DeformationGenerator -> F -> Kinematics -> invariants + energy + dW/dI
%   -> Normalizer -> MaterialDataset -> Trainer(model, EnergyStressLoss) -> trained model
\fbox{\parbox{0.9\textwidth}{\centering\vspace{3cm}
  \textbf{[PLACEHOLDER]} \\[0.5em]
  Training pipeline flowchart: deformation sampling $\to$ invariant computation
  $\to$ normalization $\to$ dual-objective training with energy and autograd
  stress.
\vspace{1cm}}}
\caption{Overview of the data generation and training pipeline.}
\label{fig:training_pipeline}
\end{figure}
